\section{Methodology}
\label{methodology}

\subsection{Framework Overview}
The proposed method treats requirement conflict detection as a consistency
verification workflow. Given a set of natural language requirements, the workflow
extracts candidate constraints, compares them against requirement knowledge, detects
potential conflicts, generates explanations, and requests human refinement when
automated evidence is insufficient.

\subsection{Multi-Agent Conflict Analysis Workflow}

\subsubsection{Analyst Agent}
The Analyst Agent decomposes each requirement into intent, actors, actions,
conditions, constraints, and expected outcomes. It also groups requirements that are
likely to interact.

\subsubsection{Validator Agent}
The Validator Agent checks the candidate interactions against consistency criteria,
constraint rules, and known semantic conflict patterns. It produces conflict
candidates with supporting evidence.

\subsubsection{Reviewer Agent}
The Reviewer Agent examines whether a reported conflict is sufficiently grounded,
whether the explanation is understandable, and whether the case should be sent to a
human stakeholder for refinement.

\subsection{Requirement Conflict Identification}

\subsubsection{Functional Conflict}
Functional conflicts occur when two requirements prescribe incompatible system
behaviors for the same condition, input, actor, or business goal.

\subsubsection{Constraint Conflict}
Constraint conflicts occur when non-functional or business constraints cannot be
satisfied together, such as mutually incompatible limits on time, access, resource
usage, or policy compliance.

\subsubsection{Role Conflict}
Role conflicts occur when requirements assign inconsistent permissions,
responsibilities, or restrictions to the same stakeholder role or to interacting
roles.

\subsubsection{Temporal Conflict}
Temporal conflicts occur when requirements impose incompatible ordering, timing,
state transition, or lifecycle constraints.

\subsection{Knowledge-Guided Consistency Analysis}

\subsubsection{Requirement Knowledge}
Requirement knowledge captures domain concepts, stakeholder roles, requirement
attributes, and consistency criteria that are needed to interpret natural language
requirements.

\subsubsection{Constraint Rules}
Constraint rules encode checkable consistency conditions, such as mutually exclusive
permissions, incompatible numeric bounds, mandatory ordering, and dependency
constraints.

\subsubsection{Semantic Patterns}
Patterns capture common conflicts, including negation, exception, role permission,
and condition outcome mismatches.

\subsubsection{Prompt Strategies}
Prompt strategies guide each agent to produce structured evidence rather than
ungrounded judgments. Prompts should require conflict type, involved requirements,
consistency rule, explanation, confidence, and suggested human review action.

\subsection{Requirement Verification and Refinement}

\subsubsection{Conflict Detection}
The workflow detects candidate conflicts by comparing requirement pairs or groups
under the selected knowledge and constraint rules.

\subsubsection{Conflict Explanation}
Each reported conflict should include the conflicting requirement statements, the
violated consistency criterion, and a concise natural language explanation.

\subsubsection{Conflict Resolution Suggestions}
The method may suggest refinement directions, such as relaxing a constraint,
clarifying an actor, adding an exception condition, or splitting a requirement.
Suggestions are treated as decision support rather than automatic requirement
rewriting.

\subsection{Human-in-the-Loop Interaction}
Human-in-the-loop interaction is used when the agents cannot determine whether a
case is a true conflict or an acceptable domain-specific exception. Human feedback is
used to confirm conflicts, reject false positives, and refine the knowledge or
constraint rules used in later verification.
